\documentclass[10pt,a4paper]{amsart}
\usepackage[margin=19mm]{geometry}
\usepackage{amsmath,amssymb,mathtools}
\usepackage[scr=rsfs,cal=euler]{mathalfa}
\usepackage{xfrac}
\usepackage[unicode,hidelinks,bookmarks=false]{hyperref}
\newcommand{\ie}{i.e.\ }
\newcommand{\cf}{cf.\ }
\newcommand{\resp}{respectively\ }
\newcommand{\Subsection}{\subsection}
\providecommand{\nobreakdash}{\nobreak\hbox{-}}
\setlength{\emergencystretch}{2em}
\multlinegap0pt
\usepackage{amssymb}
\usepackage{enumerate}
\newtheorem{lemma}{Lemma}
\newtheorem{proposition}{Proposition}
\title{\textbf{Solutions of time fractional anomalous diffusion equations with coefficients depending on both time and space variables}}
\author{Ganbileg Bat-Ochir, Khongorzul Dorjgotov, Uuganbayar Zunderiya}
\date{Source: arXiv:2405.07285v1 (CC BY 4.0)}
\begin{document}
\maketitle

\noindent\emph{Source and licence.} \textbf{Solutions of time fractional anomalous diffusion equations with coefficients depending on both time and space variables}. Version arXiv:2405.07285v1 (CC BY 4.0), distributed by arXiv under the Creative Commons Attribution 4.0 International licence, http://creativecommons.org/licenses/by/4.0/, as recorded in the arXiv licence metadata for this exact version and shown on the abstract page https://arxiv.org/abs/2405.07285v1. Full licence notice: \texttt{licenses/arxiv-2405.07285-CC-BY-4.0.txt}.\par\smallskip

\noindent\textit{Editorial scope.} Counted unit: the article's proposition thm1 (source0.tex lines 280--296). The article's complete original proof of it is delivered with it (source0.tex lines 297--383, one \texttt{proof} environment of 87 lines). No mathematics is written, repaired or re-derived: the statement and the proof are the article's own lines, in the article's own notation, and no retained line is substituted or re-flowed.  Retained uncounted as the source-local context the proof needs: lines 89--131; lines 216--242; lines 247--250; lines 254--256; lines 267--272.  Nothing was omitted from any retained span, so the package carries no omission record.  No retained line contains a citation, so the wrapper carries no bibliography and no bibliography entry.  No retained line contains a figure environment, an image-inclusion command or drawing code, so no image is delivered and the package declares no asset dependency.  Editorial formatting only, in addition: an AMS \texttt{amsart} preamble that restates the article's own declarations and macros used by the retained lines, namely the environments \texttt{lemma}, \texttt{proposition} on separate counters, together with the standard packages those lines need.  The retained environments print in this document as Lemma 1, Proposition 1 in order of appearance, whereas the article prints the same items with the numbering of its own full document.  The complete native source of the article as distributed in the arXiv e-print for this exact version is bundled unchanged as \texttt{sources/159029/source0.tex}.
We know the following identities of H-function for $z>0$
\begin{eqnarray}
H^{m,l}_{p,q}\left[z\biggr\vert\begin{array}{c}
(A_i,\alpha_i)_{1,p}\\
(B_j,\beta_j)_{1,q}
\end{array}\right] & = & H^{l,m}_{q,p}\left[\frac{1}{z}\biggr\vert\begin{array}{c}
(1-B_j,\beta_j)_{1,q}\\
(1-A_i,\alpha_i)_{1,p}
\end{array}\right],\label{eq9}\\
H^{m,l}_{p,q}\left[z\biggr\vert\begin{array}{c}
(A_i,\alpha_i)_{1,p}\\
(B_j,\beta_j)_{1,q}
\end{array}\right] & = & k H^{m,l}_{p,q}\left[z^k\biggr\vert\begin{array}{c}
(A_i,k\alpha_i)_{1,p}\\
(B_j,k\beta_j)_{1,q}
\end{array}\right]\quad\mbox{ for } k>0,\label{eq10}\\
z^\sigma H^{m,l}_{p,q}\left[z\biggr\vert\begin{array}{c}
(A_i,\alpha_i)_{1,p}\\
(B_j,\beta_j)_{1,q}
\end{array}\right] & = & H^{m,l}_{p,q}\left[z\biggr\vert\begin{array}{c}
(A_i+\sigma\alpha_i,\alpha_i)_{1,p}\\
(B_j+\sigma\beta_j,\beta_j)_{1,q}
\end{array}\right]\quad\mbox{ for any } \sigma\in\mathbb{C},\label{eq11}\\
H^{m+r,l}_{p+1,q+r}\left[z\biggr\vert\begin{array}{c}
(A_i,\alpha_i)_{1,p},(1,r)\\
\left(\frac{j}{r},1\right)_{1,r},(B_j,\beta_j)_{1,q}
\end{array}\right] & = & \frac{(2\pi)^{\frac{r-1}{2}}}{\sqrt{r}}H^{m,l}_{p,q}\left[r^r z\biggr\vert\begin{array}{c}
(A_i,\alpha_i)_{1,p}\\
(B_j,\beta_j)_{1,q}
\end{array}\right]\quad\mbox{ for } r\in\mathbb{Z}_+, l\leq p .\label{eq12}
\end{eqnarray}
Moreover the following relations hold between the special functions. The Mittag-Leffler and Wright functions can be expressed by generalized Wright functions
\begin{equation}\label{eq13}
E_{\alpha,\beta}(z) = {}_1\Psi_1\left[z\left|\begin{array}{c}
(1,1)\\
(\beta,\alpha)
\end{array}\right.
\right] \quad\mbox{and}\quad
\Psi\left(z;\alpha,\beta\right) = \sum\limits_{k=1}^\infty \frac{1}{\Gamma(\alpha k+\beta)}\frac{z^k}{k!}={}_0\Psi_{1}\left[z\left|\begin{array}{c}
-\\
(\beta,\alpha)
\end{array}\right.\right].
\end{equation}

\begin{lemma}\label{le5}
Let $\nu=\sum_{j=1}^{q}\beta_j-\sum_{i=1}^{p}\alpha_i> 0,$ $\mu=\sum\limits_{j=1}^{m}\beta_j-\sum\limits_{j=m+1}^{q}\beta_j-\sum\limits_{i=1}^{p}\alpha_i>0$, $\alpha\in \mathbb{R}_+$ and $a\in \mathbb{R}\setminus\{0\}$.  Then the following equalities hold.

(1) If $a>0$, then
\begin{multline*}
\frac{d^\alpha}{d z^\alpha}H_{p,q}^{m,0}\left[az^{-\alpha_p}\biggr\vert\begin{array}{c}
(A_i,\alpha_i)_{1,p-1}, (1, \alpha_p)\\
(B_j,\beta_j)_{1,q}
\end{array}\right]\\
=z^{-\alpha}H_{p,q}^{m,0}\left[az^{-\alpha_p}\biggr\vert\begin{array}{c}
(A_i,\alpha_i)_{1,p-1}, (1-\alpha, \alpha_p)\\
(B_j,\beta_j)_{1,q}
\end{array}\right],\quad z>0.
\end{multline*}

(2) If $m\ge 1,$ then
\begin{multline*}
\left(\frac{\beta_1}{\alpha_p}z\frac{d}{dz}+B_1\right)H_{p,q}^{m,0}\left[az^{-\alpha_p}\biggr\vert\begin{array}{c}
(A_i,\alpha_i)_{1,p-1}, (1, \alpha_p)\\
(B_j,\beta_j)_{1,q}
\end{array}\right]\\
=H_{p,q}^{m,0}\left[az^{-\alpha_p}\biggr\vert\begin{array}{c}
(A_i,\alpha_i)_{1,p-1}, (1, \alpha_p)\\
(B_1+1,\beta_1),(B_j,\beta_j)_{2,q}
\end{array}\right].
\end{multline*}
\end{lemma}

\begin{equation}\label{a1}
\frac{d^\alpha}{dz^\alpha}y(z)=
  z^m\left(a_nz^n\frac{d^n}{dz^n}y+a_{n-1}z^{n-1}\frac{d^{n-1}}{dz^{n-1}}y+\cdots+ a_1z\frac{d}{dz}y +a_0y\right), \quad z>0,
\end{equation} 

\begin{equation}\label{char}
    \widetilde  P(s)=a_n\prod\limits_{j=0}^{n-1}(s-j)+ a_{n-1}\prod\limits_{j=0}^{n-2}(s-j)+ \cdots+ a_1s +a_0
\end{equation}

\begin{lemma}\label{lem1}
Let $a\in\mathbb{R}_{+}, m\in \mathbb{N}$ and $b\in\mathbb{C}$. Then
\begin{equation}\label{lemeq}
    \frac{1}{\Gamma(1+ab+m)}\prod_{i=1}^{m}\Gamma\left(\frac{i}{a}+b+1\right)=\frac{1}{a^m\Gamma(1+ab)}\prod_{i=1}^m\Gamma\left(\frac{i}{a}+b\right).
\end{equation}
\end{lemma}

\begin{proposition}\label{thm1}
The equation given in (\ref{a1}) has the following solutions.
\begin{enumerate}
    \item For $\alpha<n,$ the solution is expressed as
    \begin{equation}\label{uuu2}
    y(z)=c_1H_{1,m+n}^{m+n,0} \left[ \frac{z^{-(\alpha+m)}}{a_n(\alpha+m)^{m+n}} \left| \begin{matrix}
( 1 , \alpha+m )  \\
\left( -\frac{s_j}{\alpha+m} , 1 \right)_{1,n}, \left(\frac{j}{\alpha+m} , 1  \right)_{1,m}   \end{matrix}\right.\right];  
    \end{equation}
    \item For $\alpha>n,$ the solution is expressed as
    \begin{equation}\label{uuu3}
        y(z)=\sum_{k=1}^{[\alpha]+1} c_kz^{\alpha-k} {}_{m+n+1}\Psi_1 \left[a_{n}(\alpha+m)^{m+n}z^{\alpha+m}\left|\begin{matrix}\left( \frac{\alpha-k-s_i}{\alpha+m} , 1 \right)_{1,n}, \left(\frac{\alpha-k+i}{\alpha+m},1\right)_{1,m} , ( 1 , 1 )\\ 
( 1+\alpha-k , \alpha+m )  \end{matrix}\right.\right].
    \end{equation}
Here  $s_1,\dots,s_n$ are the roots of the characteristic polynomial (\ref{char}). 
\end{enumerate}
\end{proposition}

\begin{proof}
Because the convergence
conditions of H-functions and generalized wright functions holds the functions in (\ref{uuu2}) and (\ref{uuu3}), it is sufficient to show that the functions satisfy the equation in (\ref{a1}).
Firstly, we rewrite the right hand side of the equation (\ref{a1}) in the factorized form
\begin{align}\label{uuu5}
    & z^{m}\left(a_nz^n\frac{d^n}{dz^n}+a_{n-1}z^{n-1}\frac{d^{n-1}}{dz^{n-1}}+\cdots+ a_1z\frac{d}{dz} +a_0\right)\nonumber\\
    &~~~~\left( c_1H_{1,m+n}^{m+n,0} \left[ \frac{z^{-(\alpha+m)}}{a_n(\alpha+m)^{m+n}} \left| \begin{matrix}
( 1 , \alpha+m )  \\
 \left( -\frac{s_j}{\alpha+m} , 1 \right)_{1,n}, & \left(\frac{j}{\alpha+m} , 1  \right)_{1,m} &  \end{matrix}\right.\right]\right)\nonumber\\ 
    &~=c_1a_n(\alpha+m)^n z^m\left(\frac{1}{\alpha+m}z\frac{d}{dz}-\frac{s_n}{\alpha+m}\right)\cdots\left(\frac{1}{\alpha+m}z\frac{d}{dz}-\frac{s_1}{\alpha+m}\right)\nonumber\\
    & ~~~~\left( H_{1,m+n}^{m+n,0} \left[ \frac{z^{-(\alpha+m)}}{a_n(\alpha+m)^{m+n}} \left| \begin{matrix}
( 1 , \alpha+m )  \\
 \left( -\frac{s_j}{\alpha+m} , 1 \right)_{1,n}, & \left(\frac{j}{\alpha+m} , 1  \right)_{1,m} &  \end{matrix}\right.\right]\right).
\end{align}
From Lemma~\ref{le5}, we get the following identity for $k=1,\ldots, n$
\begin{align*}
& \left(\frac{1}{\alpha+m}z\frac{d}{dz}-\frac{s_k}{\alpha+m}\right)\left(H_{1,m+n}^{m+n,0} \left[ \frac{z^{-(\alpha+m)}}{a_n(\alpha+m)^{m+n}} \left| \begin{matrix}
( 1 , \alpha+m )  \\
\left( -\frac{s_j}{\alpha+m} , 1 \right)_{1,n}& \left(\frac{j}{\alpha+m} , 1  \right)_{1,m} &  \end{matrix}\right.\right]\right)\\
&=H_{1,m+n}^{m+n,0} \left[ \frac{z^{-(\alpha+m)}}{a_n(\alpha+m)^{m+n}} \left| \begin{matrix}
( 1 , \alpha+m )  \\
\left( -\frac{s_j}{\alpha+m} , 1 \right)_{1,k-1}, \left(1 -\frac{s_k}{\alpha+m} , 1 \right), \left( -\frac{s_j}{\alpha+m} , 1 \right)_{k+1,n}, & \left(\frac{j}{\alpha+m} , 1  \right)_{1,m} &  \end{matrix}\right.\right].
\end{align*}
Then applying the above identity in (\ref{uuu5}), the right hand side of the equation becomes
\begin{align*}
& c_1a_nz^m(\alpha+m)^n H_{1,m+n}^{m+n,0} \left[ \frac{z^{-(\alpha+m)}}{a_n(\alpha+m)^{m+n}} \left| \begin{matrix}
( 1 , \alpha+m )  \\
 ( 1-\frac{s_j}{\alpha+m} , 1 )_{1,n} & (\frac{j}{\alpha+m} , 1  )_{1,m} &  \end{matrix}\right.\right]. 
\end{align*}
Now, let us begin the calculation of left hand side of the equation. By applying the Lemma~\ref{le5}, we get
\begin{equation*}
 \frac{d^{\alpha}}{dz^{\alpha}}y=c_1z^{-\alpha}H_{1,m+n}^{m+n,0} \left[ \frac{z^{-(\alpha+m)}}{a_n(\alpha+m)^{m+n}} \left| \begin{matrix}
( 1-\alpha , \alpha+m )\\( -\frac{s_j}{\alpha+m} , 1 )_{1,n}, (\frac{j}{\alpha+m} , 1  )_{1,m}  \end{matrix}\right.\right].   
\end{equation*}
Then applying (\ref{eq11}),  the left hand side is
$$c_1a_nz^m(\alpha+n)^{m+n}H_{1,m+n}^{m+n,0} \left[ \frac{z^{-(\alpha+m)}}{a_n(\alpha+m)^{m+n}} \left| \begin{matrix}
( 1+m , \alpha+m )\\( 1-\frac{s_j}{\alpha+m} , 1 )_{1,n}, (1+\frac{j}{\alpha+m} , 1  )_{1,m}  \end{matrix}\right.\right].$$
Hence, the left hand side becomes
\begin{equation*}
\frac{d^{\alpha}}{dz^{\alpha}}y=
c_1a_nz^m(\alpha+n)^{m+n}\frac{1}{2\pi i}\int_{L}{\frac{\prod\limits_{j=1}^n\Gamma\left( 1-\frac{s_j}{\alpha+m}-s\right)\prod\limits_{j=1}^{m} \Gamma\left(1+\frac{j}{\alpha+m}-s\right)}{\Gamma\left(1+m-(\alpha+m)s\right)}\cdot\left(\frac{z^{-(\alpha+m)}}{a_n(\alpha+m)^{m+n}}\right)^{s}\,d} s
\end{equation*}
By substituting $a=\alpha+m,$ $b=-s$ in Lemma~\ref{lem1}, we see that
\begin{align*}
&\frac{\prod\limits_{j=1}^{m} \Gamma\left(1+\frac{j}{\alpha+m}-s\right)}{\Gamma(1+m-(\alpha+m)s)}= \frac{1}{(\alpha+m)^m}\cdot\frac{\prod\limits_{j=1}^{m}\Gamma\left(\frac{j}{\alpha+m}-s\right)}{\Gamma(1-(\alpha+m)s)},
\end{align*}
which concludes the right and left hand sides of the equations are equal.

The second assertion of the theorem consists of $[\alpha]+1$ summands. Since the equation is linear, it is sufficient to show that it holds for single $k$th summand. For the fixed summand, we  will take the similar steps as in the previous case. From Lemma~\ref{le5}, we get the following identity for $l=1,\ldots, n$
\begin{align*}
& \left(\frac{1}{\alpha}z\frac{d}{dz}-\frac{s_l}{\alpha}\right)\left(   z^{\alpha-k}{}_{m+n+1}\Psi_1\left[a_{n}(\alpha+m)^{m+n}z^{\alpha+m}\left|\begin{matrix} ( \frac{\alpha-k-s_i}{\alpha+m} , 1 )_{1,n}, (\frac{\alpha-k+i}{\alpha+m},1)_{1,m} , ( 1 , 1 )\\ 
( 1+\alpha-k , \alpha+m )  \end{matrix}\right.\right] \right)=\frac{\alpha+m}{\alpha}z^{\alpha-k}\\
&~ \times {}_{m+n+1}\Psi_1 \left[a_{n}(\alpha+m)^{m+n}z^{\alpha+m}\left|\begin{matrix}( \frac{\alpha-k-s_i}{\alpha+m} , 1 )_{1,l-1}, ( \frac{\alpha-k-s_l}{\alpha+m}+1 , 1 ), ( \frac{\alpha-k-s_i}{\alpha+m} , 1 )_{l+1,n}, (\frac{\alpha-k+i}{\alpha+m},1)_{1,m} , ( 1 , 1 )\\ 
( 1+\alpha-k , \alpha+m )  \end{matrix}\right.\right],
   \end{align*}
the right hand side of the equation becomes
\begin{align*}
&a_n \alpha^{n}z^m\left(\frac{1}{\alpha}z\frac{d}{dz}-\frac{s_n}{\alpha}\right) \cdots\left(\frac{1}{\alpha}z\frac{d}{dz}-\frac{s_2}{\alpha}\right)\left(\frac{1}{\alpha}z\frac{d}{dz}-\frac{s_1}{\alpha}\right)\\
&~~~~\left(   z^{\alpha-k}{}_{m+n+1}\Psi_1 \left[a_{n}(\alpha+m)^{m+n}z^{\alpha+m}\left|\begin{matrix}( \frac{\alpha-k-s_i}{\alpha+m} , 1 )_{1,n}, (\frac{\alpha-k+i}{\alpha+m},1)_{1,m} , ( 1 , 1 )\\ 
( 1+\alpha-k , \alpha+m )  \end{matrix}\right.\right] \right)\\
& =a_n (\alpha+m)^{n} z^{\alpha+m-k}{}_{m+n+1}\Psi_1 \left[a_{n}(\alpha+m)^{m+n}z^{\alpha+m}\left|\begin{matrix}( \frac{\alpha-k-s_i}{\alpha+m}+1 , 1 )_{1,n}, (\frac{\alpha-k+i}{\alpha+m},1)_{1,m} , ( 1 , 1 )\\ 
( 1+\alpha-k , \alpha+m )  \end{matrix}\right.\right].
   \end{align*}
   The left hand side of the equation is calculated as following
    \begin{align*}
   &\frac{d^{\alpha}}{dz
^{\alpha}}z^{\alpha-k} {}_{m+n+1}\Psi_1
\left[a_{n}(\alpha+m)^{m+n}z^{\alpha+m}\left|\begin{matrix}(
\frac{\alpha-k-s_i}{\alpha+m} ,1)_{1,n},(\frac{\alpha-k+i}{\alpha+m},1)_{1,m},( 1 , 1 )\\ 
( 1+\alpha-k , \alpha+m ) \end{matrix}\right.\right]\\
&~ =b_0a_n(\alpha+m)^{m+n}z^{\alpha+m-k}\\
&~~~~\times{}_{m+n+2}\Psi_2
\left[a_{n}(\alpha+m)^{m+n}z^{\alpha+m}\left|\begin{matrix}( 1 , 1) , (
\frac{\alpha-k-s_i}{\alpha+m}+1 ,1)_{1,n},(\frac{\alpha-k+i}{\alpha+m}+1,1)_{1,m},  ( 2 , 1 )\\ 
( 2 , 1),( 1+\alpha-k+m , \alpha+m ) \end{matrix}\right.\right]\\
&~ =b_0a_n(\alpha+m)^{m+n}z^{\alpha+m-k}\\
&~~~~\times {}_{m+n+1}\Psi_1
\left[a_{n}(\alpha+m)^{m+n}z^{\alpha+m}\left|\begin{matrix} (
\frac{\alpha-k-s_i}{\alpha+m}+1 ,1)_{1,n},(\frac{\alpha-k+i}{\alpha+m}+1,1)_{1,m},  ( 1 , 1 )\\ ( 1+\alpha-k+m , \alpha+m ) \end{matrix}\right.\right]\\
&~ =b_0a_n(\alpha+m)^{m+n}z^{\alpha+m-k}\\
&~~~~\times\sum \limits_{j=0}^{\infty}\frac{ \prod\limits_{i=1}^{n}\Gamma
\left( \frac{\alpha-k-s_i}{\alpha+m}+1+j\right)  \prod\limits_{i=1}^{m}\Gamma  \left(
\frac{\alpha-k+i}{\alpha+m}+1+j\right) \Gamma(1+j)  }{\Gamma(1+\alpha-k+m+(\alpha+m)j)}\cdot\frac{
\left(a_{n}(\alpha+m)^{m+n}z^{\alpha+m}\right)^{j}}{j!}
\end{align*}
which equals to the right hand side of the equation via Lemma~\ref{lem1} with $a=\alpha+m,~b=\frac{\alpha-k}{\alpha+m}+j$.
\end{proof}

\end{document}
